\documentclass[a4paper,11pt]{article}
\usepackage{sbpo-template}
\usepackage[brazil]{babel}
\usepackage[utf8]{inputenc}
\usepackage[T1]{fontenc}
\usepackage{amsmath,amssymb}
\usepackage{url}

\usepackage[square]{natbib}
\usepackage{indentfirst}

\title{TÍTULO DO ARTIGO}

\begin{document}

\maketitle

\author{
\name{Primeiro Autor}
\institute{Filiação}
\iaddress{Endereço da Instituição}
\email{e-mail}
}

\author{
\name{Segundo Autor}
\institute{Filiação}
\iaddress{Endereço da Instituição}
\email {e-mail}
}

\vspace{8mm}
\begin{resumo}
Este modelo resume as normas de formatação para os trabalhos completos a serem publicados como produção bibliográfica da disciplina de GED-16. Título, filiação, resumo e palavras-chave devem repetir fielmente o que foi informado quando o autor apresentou o Projeto Final na disciplina.
O Resumo deve ter no máximo 150 palavras.
 \end{resumo}

\bigskip
\begin{palchaves}
Primeira, Segunda, Terceira.

\end{palchaves}


\vspace{8mm}

\begin{abstract}
Aqui deverá vir o Abstract em inglês
e as \textit{keywords}, traduzindo fielmente o Resumo e as palavras-chave.
The Abstract must not exceed 150 words.
\end{abstract}

\bigskip
\begin{keywords}
First keyword. Second keyword. Last keyword.

\end{keywords}


\newpage

\section*{ORIENTAÇÕES GERAIS}

1. Os nomes que não sejam informados nesse formulário não aparecerão entre os autores na publicação da coletânea de artigos da disciplina.


2. Este manuscrito será disponibilizado para o exame pelos revisores (docentes da disciplina). Por isso, devem ser entregues tanto o arquivo .tex quanto uma versão salva em .pdf.

3. Uma vez avaliado o artigo, os autores serão chamados a encaminhar \textbf{versão final} com as correções requeridas pelas docentes.

4. As páginas deste texto não devem vir numeradas, tanto no caso de arquivo enviado na data de entrega do projeto, quanto no caso do arquivo com a versão final do artigo corrigido.
A numeração será feita posteriormente para o conjunto de todos os artigos.

5. Cabeçalhos e rodapés devem ser deixados em branco.

\vspace{0.2cm}

\section*{ Instruções de Formatação}
Os trabalhos completos devem ter \textbf{no máximo 15 páginas}, incluídos neste limite: a primeira página com resumo, texto, tabelas, gráficos, agradecimentos e referências.

Os textos devem utilizar páginas de tamanho \textbf{A4} (29,7 x 21,0 cm) com \textbf{margem superior de 3,3 cm, inferior de 2,5 cm e laterais de 2,9 cm}.
 Devem ser escritos em coluna única, com fonte \textbf{\textit{Times New Roman} 11}.

\vspace{1cm}

\noindent \textbf{O artigo final deve ser organizado em introdução, questão de pesquisa, objetivo(s), metodologia, resultados, discussão e conclusão.}

\section{Introdução}
Descrever o problema abordado, justificando a relevância do projeto em sua área.

Em seguida, apresentem a questão de pesquisa e o(s) objetivo(s) do estudo.

\subsection{Questão de Pesquisa}
Coloque a Questão de Pesquisa proposta

\subsection{Objetivos}
Enfatizar o objetivo geral \textbf{e/ou }específicos a que este trabalho se destina.

\vspace{0.2cm}

\textbf{Objetivo Geral:}

\textbf{Objetivos Específicos:}


\section{Metodologia}
Na seção de Metodologia, descrevam a base de dados, as variáveis analisadas, o período em que as mesmas foram coletadas etc. Em seguida, descreva os métodos estatísticos utilizados, explicando a finalidade de cada um.

\section{Resultados}
Descreva os resultados finais obtidos no desenvolvimento do seu trabalho. Você pode utilizar todos os recursos que possam facilitar a compreensão dos resultados, como tabelas, gráficos, fotos ou ilustrações.

\subsection{Discussão}
Faça a avaliação dos seus resultados, identifique pontos positivos e negativos do seu trabalho, mostre qual foi o seu raciocínio e como os dados encontrados colaboraram na evolução do seu pensamento e na formulação de suas conclusões.

\vspace{0.2cm}
\noindent \textbf{OBS.:} As seções de Resultados e Discussão, se quiser, pode vir conjuntamente em uma única seção de Resultados e Discussão, na qual devem ser apresentados e interpretados os principais achados em relação à questão de pesquisa.

\section{Considerações Finais}
Baseado em seus resultados e em sua discussão elabore a sua conclusão. Ela deve ser uma resposta ao que foi proposto no objetivo do trabalho.

\subsection{Limitações}
Indiquem as limitações do estudo e sugiram formas de aprofundá-lo em trabalhos futuros.

\vspace{1cm}
\noindent \textbf{Por fim, incluam as referências citadas no texto. Alguns exemplos estão dispostos a seguir:}

\vspace{0.2cm}
\bibliographystyle{sbpo}
\bibliography{exemplo-latex}

\vspace{1cm}
\textbf{Estilo das Citações}

As citações no texto devem estar entre colchetes e conter  os últimos sobrenomes dos autores~\citep{silva:99}, no caso de um ou dois autores, e o último sobrenome seguido de "et al." no caso de mais de dois autores, seguidos do \textbf{ano da publicação}, como por exemplo,~\citep{anna:06},~\citep{gates:03}, ~\citep{smith:02}, ~\citep{silva:99}, ~\citep{pele:04}, ~\citep{web:16}.
As referências no final do texto devem estar em ordem alfabética do último sobrenome do primeiro autor.

\end{document}
